\sectiontitle{一、单项选择题}{10}{20}

% B01 S01:47
\mc{关于自动驾驶感知任务，下列对应关系正确的是：}
{分类：持续估计同一目标的运动状态}
{跟踪：只判断一张图像中是否存在汽车}
{目标检测：只输出整张图像的类别，不输出位置}
{语义分割：为像素或点分配语义类别}
{D}{语义分割给出像素或点级的类别。检测通常需要位置，跟踪需要在时间上关联目标。}

% B02 S05:28-31
\mc{半径离群点滤波在点云处理中，主要依据什么判断某点是否可能为孤立点？}
{该点给定半径邻域内的邻居数量}
{该点离传感器最远就无条件删除}
{该点反射强度最大就无条件删除}
{该点是否位于某一坐标轴的正半轴}
{A}{给定半径邻域内邻点过少时，该点可能为孤立点。阈值需要结合点云密度设置。}

% B03 S03:17-19
\mc{采用课件的 OpenDRIVE 道路坐标 $(s,t)$ 时，$t$ 表示：}
{传感器采样时间}
{沿道路参考线的累计纵向位置}
{相对道路参考线的横向位置}
{相对道路参考线的航向角}
{C}{此处 $t$ 是道路坐标符号，表示横向位置，不能与 S-T 图中的时间混淆。}

% B04 S10:9-10
\mc{城市高楼附近，GNSS 接收到了建筑物反射后的卫星信号。这主要对应哪类问题？}
{相机曝光不足}
{惯性测量误差的积分漂移}
{多路径效应}
{车体与雷达之间的外参误差}
{C}{反射使传播路径与直接路径不同，会影响距离估计，属于多路径问题。}

\exampagebreak

% B05 S06:22-24;S07:19-22
\mc{为什么行人轨迹预测应考虑行人与行人的交互？}
{行人未来位置完全由当前速度唯一决定}
{避让、同行和群组行为会改变未来运动}
{只要加入交互，预测就一定没有误差}
{交互信息可以替代所有位置观测}
{B}{同行、群组与避让等行为使运动具有交互性。交互建模不能消除全部不确定性。}

% B06 H3:25-27
\mc{在横轴为时间 $t$、纵轴为路径累计距离 $s$ 的 S-T 图中，曲线斜率 $\mathrm{d}s/\mathrm{d}t$ 表示：}
{沿路径的加速度}
{沿路径的累计距离}
{沿路径的加加速度}
{沿路径的速度}
{D}{$s$ 的单位为 m，$t$ 的单位为 s，因此斜率的单位为 m/s。}

% B07 S02:26-28;H4:27-29
\mc{按作业四的 LQR 参数约定，增大 $q_0$ 最直接改变的是：}
{代价函数对横向位置误差的惩罚权重}
{代价函数对航向误差的惩罚权重}
{参考路径的道路限速}
{代价函数对控制输入的惩罚权重}
{A}{$q_0$ 对应横向位置误差权重。实际响应还取决于模型、其他权重与执行约束，不能只凭 $q_0$ 推断性能一定提高。}

% B08 S09:16-17
\mc{DQN 训练中，经验回放保存的核心内容通常是：}
{仅每回合结束时的总耗时}
{状态、动作、奖励、下一状态等转移样本}
{只保存当前网络参数的梯度}
{只保存动作名称而不保存状态}
{B}{核心转移为 $(s,a,r,s')$。实现中还需记录终止信息等，以正确构造训练目标。}

% B09 H1:20-23
\mc{按作业一给出的命令，循环回放离线数据包应使用：}
{\code{cyber_recorder record -a}}
{\code{cyber_monitor}}
{\code{cyber_recorder play -f demo_3.5.record -l}}
{\code{bash scripts/bootstrap_neo.sh}}
{C}{该命令在作业中用于指定数据包并循环回放。题目依据课程实验环境，不要求在其他版本上执行。}

% B10 S02:18
\mc{PID 微分项对高频测量噪声较敏感的主要原因是：}
{微分项把所有误差累积后求平均}
{微分项对快速变化分量总是抑制最强}
{微分项只保留恒定误差而删除变化分量}
{差分或微分会放大快速变化的测量分量}
{D}{变化率运算会增强高频噪声的影响，因此实际控制中常需要滤波并谨慎选择微分增益。}

\exampagebreak
\sectiontitle{二、填空题}{10}{20}
\textbf{说明：}每题一个空，每空 2 分。数值题应按给定公式和符号约定作答。

% B11 S10:15
\fillq{实时激光点云与已有地图点云之间可使用迭代最近点算法配准。该算法的英文缩写为 \blankline。}
{ICP}{ICP 的全称为 Iterative Closest Point。}

% B12 S05:38-44
\fillq{KD-tree 通过沿坐标维度递归划分 \blankline，缩小近邻搜索需要访问的区域。}
{空间（或数据空间）}{不能把 KD-tree 描述为仅按点在文件中的顺序分组。}

% B13 H3:25-27
\fillq{S-T 图中，横轴单位为 s，纵轴单位为 m。曲线斜率的单位为 \blankline。}
{m/s}{距离对时间的一阶导数对应速度。}

% B14 S10:6,22
\fillq{在定位融合算法中，KF 的中文名称为 \blankline。}
{卡尔曼滤波}{也接受“卡尔曼滤波器”。EKF 和 UKF 分别为其不同扩展形式。}

% B15 S02:27-28
\fillq{LQR 的代价写成 $\int(x^\mathsf{T}Qx+u^\mathsf{T}Ru)\,\mathrm{d}t$ 时，矩阵 \blankline[1.5cm] 用于惩罚控制输入。}
{$R$}{$Q$ 惩罚状态偏差，$R$ 惩罚控制输入。}

% B16 S06:17
\fillq{一维恒加速度预测采用 $x=x_0+v_0t+\frac12at^2$。已知 $x_0=1$ m、$v_0=2$ m/s、$a=2$ m/s$^2$、$t=3$ s，则 $x=$\blankline。}
{$16$ m}{代入得 $x=1+2\times3+\frac12\times2\times3^2=1+6+9=16$ m。}

% B17 S02:15-19
\fillq{使用 $u=K_pe+K_iI+K_dD$。已知 $K_p=2$、$K_i=0.5$、$K_d=0.1$，
$e=2$、$I=3$、$D=-1$，则 $u=$\blankline。各量采用相容的归一化单位。}
{$5.4$}{比例、积分、微分项分别为 $4$、$1.5$、$-0.1$，因此 $u=4+1.5-0.1=5.4$。}

% B18 S09:14-16
\fillq{Q-learning 的某次转移到达终止状态。规定终止后未来回报为 0，即时奖励为 $-2$，则该转移的一步训练目标为 \blankline。}
{$-2$}{终止转移不再加入下一状态的自举价值，因此目标等于即时奖励 $r=-2$。}

% B19 S07:23-24
\fillq{单条预测轨迹的 ADE 定义为各预测时刻位置误差欧氏距离的平均值。
三个时刻的距离误差为 $1$、$2$、$3$ m，则 ADE 为 \blankline。}
{$2$ m}{$\mathrm{ADE}=(1+2+3)/3=2$ m。不是三个距离的平方平均。}

% B20 S04:14-15
\fillq{图像直线为 $y=-2x+140$。在 $y=100$ 的水平线上，交点的 $x$ 坐标为 \blankline。坐标单位为 pixel。}
{$20$ pixel}{由 $x=(100-140)/(-2)=20$ 得到。}

\exampagebreak
\sectiontitle{三、问答题}{3}{30}

% B21 S05:28-31,38-44;S03:7-9;S06:11-17
\essay{系统已经采集激光点云，并获得高精度地图和车辆历史轨迹。请说明这些信息如何支持后续预测。
\begin{subquestions}
\item 说明点云离群点处理和近邻搜索的作用，并指出滤波过强的一项风险。（4 分）
\item 说明高精度地图能向车道序列预测提供哪些信息。（3 分）
\item 说明为什么车辆历史轨迹和当前速度仍有必要，不能只根据地图预测邻车未来位置。（3 分）
\end{subquestions}}
{\textbf{参考要点与评分：}
\begin{subquestions}
\item 离群点处理降低孤立噪声干扰，得 1 分；半径法或统计法的依据说明正确得 1 分。
近邻搜索支持局部结构分析、配准或聚类等，任一用途得 1 分。
阈值过强可能删除稀疏远处目标的有效点，或破坏小目标，得 1 分。
\item 车道中心及几何走向，得 1 分；车道连接与路口拓扑，得 1 分；
边界、可通行关系或道路语义约束，任一项得 1 分。
\item 地图提供静态可行结构，得 1 分；历史位置、速度或加速度反映当前运动趋势，得 1 分；
同一路口可能存在转弯、直行等多种意图，必须结合动态状态或交互信息，得 1 分。
\end{subquestions}}

% B22 H3:2-12,19-28
\essay{在速度规划 S-T 图中，横轴为 $t$，纵轴为 $s$。
某候选轨迹在 $0\le t\le3$ 内为 $s(t)=5t$。
障碍占据区域为 $1\le t\le2$、$4\le s\le6$。
该区域已经包含本题所需的车辆外形和安全裕度。
\begin{subquestions}
\item 说明 S-T 曲线一阶导数与二阶导数的物理意义。（2 分）
\item 判断该候选是否与占据区域相交，给出一个交点或相交时间区间。（3 分）
\item 按作业三设计“加速超车”和“减速让行”的速度数据对比，列出必须固定、采集和比较的内容。（5 分）
\end{subquestions}}
{\textbf{参考要点与评分：}
\begin{subquestions}
\item 一阶导数为速度，得 1 分；二阶导数为加速度，得 1 分。
\item $4\le5t\le6$ 给出 $0.8\le t\le1.2$，得 1 分；
与障碍时间区间 $[1,2]$ 求交得到 $[1,1.2]$，得 1 分；
因此相交，在本题约束下该候选不可行，得 1 分。只要举出 $(t,s)=(1,5)$ 并说明属于双方，也可得满分。
\item 固定场景初态、地图和参数版本，得 1 分；
记录场景名称、参数值和时间戳，得 1 分；
保存 $s(t)$、速度和加速度数据或相应曲线，得 1 分；
结合障碍占据区域比较避碰与让行行为，得 1 分；
比较限速、加速度限制和变化平顺性，得 1 分。仅按不同曲线截图判断“更安全”而无时间对应关系，不得相应证据分。
\end{subquestions}}

% B23 S02:16-18,26-33;H4:5,27-29
\essay{某车需要纵向跟踪目标速度，并在弯道保持参考轨迹。
\begin{subquestions}
\item 说明 PID 中 P、I、D 的主要作用及一个调参风险。（3 分）
\item 说明 LQR 的状态误差、$Q$ 和 $R$ 分别表达什么。（3 分）
\item 说明 MPC 如何处理未来预测与执行约束，以及计算时间为何也是设计条件。（4 分）
\end{subquestions}}
{\textbf{参考要点与评分：}
\begin{subquestions}
\item P 对当前误差响应，I 累积误差，D 响应误差变化率，三者正确得 2 分。
增益过大引发振荡、积分累积或微分放大噪声，任一合理风险得 1 分。
\item 可用横向误差、航向误差及其变化率构造状态，得 1 分；
$Q$ 惩罚状态误差，得 1 分；$R$ 惩罚控制输入，得 1 分。
\item 用车辆模型预测未来状态并优化一串控制量，得 1 分；
将转角、加速度等限制加入优化，得 1 分；
每周期只执行第一项并重新优化，得 1 分；
必须在控制周期内完成求解，超时会破坏实时执行条件，得 1 分。
\end{subquestions}
不把“预测范围更长”或“求解时间更多”视为性能必然提高的充分条件。}

\exampagebreak
\sectiontitle{四、程序题}{2}{30}

% B24 S04:13-15
\programq{编写 \code{merge_lanes(segments, y_bottom, y_top)}，将候选线段合并为左右两条直线。

\textbf{输入：}候选线段序列，每项为有限数 \code{(x1,y1,x2,y2)}。图像中 $x$ 向右、$y$ 向下。
可以没有线段，且 \code{y_bottom > y_top}。本题采用斜率符号分组的简化直线路段模型。

\begin{subquestions}
\item 跳过 $x_1=x_2$ 的线段。对其余线段计算
$k=(y_2-y_1)/(x_2-x_1)$、$b=y_1-kx_1$。
丢弃 $|k|<0.5$ 的线段，$k<0$ 分为左组，$k>0$ 分为右组。（5 分）
\item 对每组的 $k$ 和 $b$ 分别取算术平均。采用 $x=(y-b)/k$，
计算在 \code{y_bottom} 与 \code{y_top} 上的端点。（5 分）
\item 返回字典，键为 \code{"left"} 和 \code{"right"}。
非空组对应 \code{(x_bottom,y_bottom,x_top,y_top)}，空组对应 \code{None}。
坐标保持浮点数，不取整、不裁剪。写出以下样例的输出。（5 分）
\end{subquestions}
\codefile{snippets/B01.py}}
{\textbf{参考程序：}
\codefile{programs/B24.py}
\textbf{样例过程：}左组斜率均为 $-2$，截距为 140、144，故平均为 $(-2,142)$。
右组斜率均为 2，截距为 $-60$、$-56$，故平均为 $(2,-58)$。
第五条线段斜率为 $0.1$，第六条为垂直线，均被跳过。
\[
\begin{aligned}
x_{\mathrm{left,bottom}}&=(100-142)/(-2)=21,\\
x_{\mathrm{left,top}}&=(60-142)/(-2)=41,\\
x_{\mathrm{right,bottom}}&=(100+58)/2=79,\\
x_{\mathrm{right,top}}&=(60+58)/2=59.
\end{aligned}
\]
返回：
\codefile{snippets/B02.py}
\textbf{评分：}垂直线检查 1 分，斜率和截距 2 分，阈值与左右分组 2 分；
两组分别平均 3 分，交点公式 2 分；
缺失组返回 \code{None} 2 分，字典结构 1 分，样例结果 2 分。
不能将所有线段先混合平均后再分左右。}

\newpage

% B25 H3:7-12,22-28
\programq{速度规划数据已整理成 \code{rows}，每条记录为 \code{(t,s,v)}。
编写 \code{analyze_speed(rows, speed_limit, accel_limit, distance_tol)}。

\textbf{输入：}至少两条记录。$t$、$s$、$v$ 均为有限数，单位分别为 s、m、m/s。
速度非负，三个限制参数非负。输入结构有效，但时间可能重复或倒序。

\begin{subquestions}
\item 对相邻记录，若 $\Delta t\le0$，抛出 \code{ValueError}，不得静默排序或跳过。
否则按 $a_i=(v_i-v_{i-1})/(t_i-t_{i-1})$ 求区间加速度。（5 分）
\item 返回最大速度、最大绝对加速度，以及所有超速记录下标和加速度超限区间下标。
分别使用 $v_i>\text{speed\_limit}$ 与 $|a_i|>\text{accel\_limit}$。区间用右端记录下标表示。（5 分）
\item 按梯形公式检查每个区间的数据一致性：
\[
\left|(s_i-s_{i-1})-\frac{v_{i-1}+v_i}{2}(t_i-t_{i-1})\right|>
\text{distance\_tol}.
\]
返回所有不一致区间下标，并写出样例的完整输出。（5 分）
\end{subquestions}
字典键固定为下列样例说明中的五个键。所有下标从 0 起算，列表按下标递增。等于限制时不算超限。
\codefile{snippets/B03.py}}
{\textbf{参考程序：}
\codefile{programs/B25.py}
\textbf{样例过程：}三个时间间隔均为 1 s，各区间的加速度均为
$(v_i-v_{i-1})/1=2$ m/s$^2$，因此最大绝对加速度为 2 m/s$^2$，
区间 1、2、3 均超过 1.5 m/s$^2$。
最大速度为 8 m/s；仅下标 3 的速度严格大于 6 m/s。

梯形公式给出的区间距离为 3、5、7 m，实际记录增量为 3、5、8 m，
差值绝对值为 0、0、1 m。因此只有区间 3 不一致，输出为：
\codefile{snippets/B04.py}
\textbf{评分：}重复或倒序时间抛异常 2 分，正确差分加速度 3 分；
两个最大值 2 分，超速下标 1 分，加速度超限下标 2 分；
梯形距离比较 2 分，区间下标及字典结构 1 分，样例完整结果 2 分。
距离一致性是本题给定的离散检查，不能据此断言记录具有完整物理真实性。}
